%%%%%%%%%%%%%%%%%%%%%%%%%%%%%%%%%%%%%%%%%%%%%%%%%%%%%%%%%%%%%%%%
\section{Outlook and Future Work}
\label{sec:discussion_outlook}

% The limitations of the thesis lead to several concrete extensions.
% The most direct next step is experimental validation.
% A physical pendulum setup, or another controlled laboratory system with contact and measurable actuator behavior, would allow the numerical reference comparisons of Chapter~\ref{chap:case_study} to be complemented by physical data.
% This would separate modeling errors from co-simulation errors and would clarify how well the contact and switching assumptions represent a real system.

% A second direction is improved state transfer for runtime model switching.
% The current \texttt{MasterPendulum} reconstructs the FEM state from rigid-body quantities when switching into the high-fidelity model.
% Future work could preserve additional reduced deformation information, use a precomputed modal representation, or initialize the FEM model through a short relaxation step before contact.
% Such extensions would reduce the difference between switched and full-FEM trajectories while keeping the computational benefit of selective high-fidelity activation.

The hybrid co-simulation algorithm could be extended by adaptive and error-controlled macro stepping.
The present implementation uses prescribed macro step sizes.
An adaptive algorithm would estimate coupling error, reject inaccurate macro steps, and repeat them with smaller step sizes.
This requires a stronger rollback contract across all components that participate in the rejected step.
It also requires clear rules for backend components that cannot restore arbitrary past states.

The dense-time representation could be made more rigorous.
The current implementation combines floating-point physical time with an integer microstep index.
Future versions could use an exact tick-based time representation for physical time and reserve the microstep index only for event ordering at the same instant.
This would reduce tolerance-dependent comparisons in long hybrid simulations and would align better with deterministic hybrid co-simulation semantics.

The algebraic-loop solver also offers room for improvement.
The current implementation solves algebraic loops in the interface form used by \syssimx{}, but it does not reuse Jacobians across iterations or time steps.
For loops with a constant or slowly changing local linearization, Jacobian reuse could reduce computational cost.
A more complete implementation of interface-Jacobian based co-simulation would also allow stronger convergence management and clearer diagnostics for difficult loop systems.

Further work should also address scalability.
Independent execution generations are currently evaluated sequentially.
Parallel execution would be useful for larger systems, especially when many components can be advanced independently within one communication step.
A systematic benchmark suite should therefore include larger component graphs, different coupling densities, and different backend costs.
This would clarify when orchestration overhead dominates and when high-fidelity model switching provides the largest benefit.

The backend integrations can be broadened as well.
FMI~3.0 concepts such as clocks and scheduled execution are relevant for sampled and event-driven co-simulation workflows.
Supporting these features would improve the treatment of hybrid and discrete-time components beyond the FMI~2.0 Co-Simulation capabilities used in this thesis.
The OpenSim and FEM integrations could also be generalized.
In particular, a reusable FEM backend wrapper would separate NGSolve-specific model code from the common co-simulation interface more clearly.

Finally, the user-facing diagnostics of the framework should be expanded.
The structural-analysis metadata, algebraic-loop blocks, event traces, switching decisions, and benchmark measurements are already available internally or through selected visualizations.
Future versions could expose this information through a more systematic reporting layer.
This would make assembled simulations easier to inspect, reproduce, and debug, which is essential for applying heterogeneous hybrid co-simulation to larger systems.